\documentclass[
    twoside,
    10pt
]{book}

%========================================%
% PAGE SETTINGS                          %
%========================================%

\usepackage[
    paperwidth = 11cm,
    paperheight = 15cm,
    margin = 0.8cm,
    footskip = 0.5cm,
    % showframe
]{geometry}

\usepackage[utf8]{inputenc}
\usepackage[T1]{fontenc}
\usepackage{helvet}
\renewcommand{\familydefault}{\sfdefault}

\usepackage[
	skip = 0pt,
	indent = 0pt
]{parskip}

\usepackage{graphicx}
\graphicspath{ {./images} }

\usepackage{tabularray}

\usepackage{textcomp}


%========================================%
% PAGE STYLE                             %
%========================================%

\usepackage{fancyhdr}

\fancypagestyle{std}{
    \renewcommand{\headrulewidth}{0pt}
	\fancyfoot[C]{-- \thepage\ --}
}


%========================================%
% CLORS                                  %
%========================================%

\usepackage{xcolor}

\definecolor{main}{RGB}{0, 60, 160}
\color{main}


%========================================%
% TIKZ                                   %
%========================================%

\usepackage{tikz}
\usepackage{tikzpagenodes}
\renewcommand{\frame}{
    \tikz[remember picture, overlay] \draw
        (current page text area.south west)
        rectangle
        (current page text area.north east);
}


%========================================%
% TABLE                                  %
%========================================%

\NewTblrEnviron{hhtable}
\SetTblrInner[hhtable]{
    width = {\textwidth},
    rows = {
        ht = 0.84cm,
        valign = m,
        abovesep = 0.12cm
    },
    row{1} = {
        ht = 1.5cm,
        valign = m,
        abovesep = 0.2cm,
        font = {\large\bfseries}
    },
    row{2} = {
        abovesep = 0.1cm
    },
    hlines, vlines
}


%========================================%
% TOC                                    %
%========================================%

\usepackage{tocloft}
\renewcommand{\contentsname}{}
\renewcommand{\cfttoctitlefont}{}
\renewcommand{\cftaftertoctitle}{}
\setlength{\cftbeforetoctitleskip}{0pt}
\setlength{\cftaftertoctitleskip}{0pt}


\begin{document}
    \frontmatter
    \pagestyle{empty}
    \begin{titlepage}
        \centering
        \includegraphics[height = 3cm]{logo.pdf} \\[0.5cm]
        \textbf{\Large A.S.D. DOJO Trieste}
        \vfill
        \textbf{Libretto personale del} \\[0.4cm]
        \textbf{\Huge SOCIO}
        \vfill
        \includegraphics[height = 1cm]{fijlkam.pdf} \\
        Associazione sportiva affiliata \\
        alla FIJLKAM e al CONI
    \end{titlepage}
    \mainmatter{}
    \pagestyle{std}

    %
    % Person data
    %
    \addcontentsline{toc}{chapter}{DATI SOCIO}
    \frame
    \begin{center}
        \textit{Associazione Sportiva Dilettantistica} \\[0.2cm]
        \textbf{\LARGE DOJO TRIESTE}
    \end{center}
    \vfill
    \begin{center}
        \begin{tikzpicture}
            \draw (0,0) rectangle (35mm,45mm) node[midway, text opacity = 0.5] {FOTO};
        \end{tikzpicture}
    \end{center}
    \vfill\vfill
    \begin{tblr}{
        width = {\textwidth},
        colspec = {l X},
        hline{2-Y} = {solid},
        row{Z} = {
            belowsep = 0.3cm
        }
    }
        COGNOME & \\
        NOME \\
        DATA DI NASCITA \\
        LUOGO DI NASCITA \\
        CODICE FISCALE \\
        RESIDENZA \\
        DATA DI EMISSIONE
    \end{tblr}


    %
    % Associations
    %
    \newcommand{\associations}{
        \begin{hhtable}{
            colspec = {c c X[c]},
            cell{1}{1} = {c=3}{c},
            cell{3-Z}{1} = {cmd=\inteval{\value{rownum}-2}}
        }
            ASSOCIAZIONI \\
            N. & ANNO SOCIALE & TIMBRO E FIRMA \\ \\ \\ \\ \\ \\ \\ \\ \\ \\ \\
        \end{hhtable}
    }

    \newpage
    \addcontentsline{toc}{chapter}{ASSOCIAZIONI}
    \associations
    
    \newpage
    \associations

    %
    % Athlete
    %
    \newcommand{\athlete}{
        \begin{hhtable}{
            colspec = {c c X[c]},
            cell{1}{1} = {c=3}{c},
            cell{3-Z}{1} = {cmd=\inteval{\value{rownum}-2}}
        }
            REGISTRO ATTIVITÀ \\
            N. & ANNO SPORTIVO & TIMBRO E FIRMA \\ \\ \\ \\ \\ \\ \\ \\ \\ \\ \\
        \end{hhtable}
    }

    \newpage
    \addcontentsline{toc}{chapter}{REGISTRO ATTIVITÀ}
    \athlete
    
    \newpage
    \athlete

    %
    % TOC
    %
    \newpage
    \tableofcontents
    \thispagestyle{empty}
\end{document}